
\chapter{CPSBench Battery Track}
\label{ch:cpsbench}

Standard forecasting benchmarks evaluate prediction accuracy in isolation:
they score point forecasts or intervals against held-out actuals without
ever closing the loop through a physical plant.  The CPSBench Battery Track
closes that loop, coupling forecasts to a battery dispatch controller and
evaluating safety on the \emph{true} (unobserved) state trajectory under
replayable telemetry faults.

% ----------------------------------------------------------------
\section{Why Standard Evaluation Is Insufficient}

A forecast benchmark that reports MAE and PICP answers only half the
question.  Consider two controllers:

\begin{itemize}
  \item Controller~A achieves PICP@90 = 92\% with narrow intervals,
    maximising economic value---but during the 8\% of uncovered steps,
    the dispatched action drives the battery SOC below the safe floor,
    triggering a BMS trip that takes the asset offline for 30 minutes.
  \item Controller~B achieves PICP@90 = 96\% with wider intervals,
    sacrificing 3\% of revenue---but never triggers a BMS trip.
\end{itemize}

Standard metrics rank~A above~B on interval sharpness and close to~B on
coverage.  Yet in operational terms, A is unsafe and B is safe.  The
missing ingredient is \emph{closed-loop evaluation on the true plant
state}, which requires:
\begin{enumerate}
  \item A plant simulator with deterministic SOC dynamics.
  \item A fault injector that degrades telemetry between the plant and the
    controller.
  \item Separate logging of truth-state and observed-state trajectories.
  \item Safety metrics computed exclusively on the truth trajectory.
\end{enumerate}

CPSBench provides all four.

% ----------------------------------------------------------------
\section{Dual-Trajectory Evaluation}

The central architectural choice in CPSBench is the maintenance of two
parallel trajectories at every step~$t$:

\begin{description}
  \item[Truth trajectory $\{z_t^{\mathrm{true}}\}$.]
    The plant simulator advances the true SOC using the
    \emph{executed} action and the true load/renewables signal.  This
    trajectory is logged but never exposed to the controller.

  \item[Observed trajectory $\{\tilde{z}_t\}$.]
    The fault injector corrupts the truth signal before delivering it to
    the controller.  The controller makes all decisions based on this
    degraded view.
\end{description}

Safety violations are defined as events where the truth-trajectory SOC
exits the safe envelope $[E_{\min}, E_{\max}]$.  This is the
\emph{true-SOC violation} (as opposed to an observed-SOC violation, which
may be masked by sensor error).

The dual-trajectory architecture ensures that the benchmark answers the
correct question: ``did the battery actually violate its SOC limits?''
rather than ``did the controller \emph{think} it violated?''

% ----------------------------------------------------------------
\section{Truth-State and Observed-State Logging}

At each step the CPSBench runner logs the following record to the episode
trace (cf.\ \texttt{reports/publication/dc3s\_step\_trace.csv}):

\begin{table}[ht]
\centering
\caption{Per-step logging schema in CPSBench.}
\label{tab:cpsbench-log-schema}
\begin{tabular}{lll}
\toprule
\textbf{Field} & \textbf{Source} & \textbf{Trajectory} \\
\midrule
\texttt{load\_observed\_mw}    & Fault injector   & Observed \\
\texttt{renewables\_observed\_mw} & Fault injector & Observed \\
\texttt{renewables\_true\_mw}  & Plant simulator  & Truth \\
\texttt{reliability\_w}        & OQE              & Observed \\
\texttt{drift\_flag}           & Page--Hinkley    & Observed \\
\texttt{inflation}             & RAC-Cert         & Computed \\
\texttt{proposed\_charge\_mw}  & Optimiser        & Decision \\
\texttt{proposed\_discharge\_mw} & Optimiser      & Decision \\
\texttt{safe\_charge\_mw}      & DC3S shield      & Decision \\
\texttt{safe\_discharge\_mw}   & DC3S shield      & Decision \\
\texttt{intervened}            & Shield           & Decision \\
\texttt{soc\_after\_mwh}       & Plant simulator  & Truth \\
\texttt{guarantee\_passed}     & Certificate      & Computed \\
\bottomrule
\end{tabular}
\end{table}

Every field is deterministic for a given (scenario, seed, controller)
triple, enabling exact reproduction of any individual step.

% ----------------------------------------------------------------
\section{Fault Taxonomy and Replayable Schedules}

CPSBench implements six fault types, each controlled by a severity parameter
and applied by a deterministic injector seeded from the episode RNG:

\begin{table}[ht]
\centering
\caption{Fault taxonomy with severity parameterisation.}
\label{tab:fault-taxonomy-cpsbench}
\begin{tabular}{llll}
\toprule
\textbf{Fault} & \textbf{Parameter} & \textbf{Sweep levels} &
  \textbf{Effect on $\tilde{z}_t$} \\
\midrule
Dropout       & Rate $p$         & 0, 0.05, 0.10, 0.20, 0.30 &
  $\tilde{z}_t = \mathrm{NaN}$ w.p.\ $p$ \\
Delay jitter  & Max delay (s)    & 0, 1, 5, 15 &
  $\tilde{z}_t = z_{t-\delta}$, $\delta \sim U[0, d_{\max}]$ \\
Spike         & Scale $\sigma$   & 0, 1, 2, 3 &
  $\tilde{z}_t = z_t \cdot (1 + \epsilon)$, $\epsilon \sim N(0, \sigma^2)$ \\
Stale sensor  & Window (steps)   & 0, 2, 4, 8 &
  $\tilde{z}_t = z_{t-s}$ for $s$ consecutive steps \\
Drift         & Rate (MW/step)   & 0, 0.5, 1.0, 2.0 &
  $\tilde{z}_t = z_t + \beta \cdot t$ \\
Reorder       & Rate $r$         & 0, 0.05, 0.15 &
  Adjacent packets swapped w.p.\ $r$ \\
\bottomrule
\end{tabular}
\end{table}

Each fault schedule is a deterministic function of the episode seed,
meaning that any researcher with the same code and seed will reproduce
exactly the same fault pattern.  The \texttt{drift\_combo} scenario
applies all six faults simultaneously at moderate severity, representing a
worst-case but field-plausible operating condition.

% ----------------------------------------------------------------
\section{Benchmark Metrics}

The CPSBench Battery Track reports three families of metrics, computed
exclusively on the truth trajectory:

\subsection{Safety Metrics}

\begin{description}
  \item[TSVR (True-SOC Violation Rate).]
    Fraction of steps where the true SOC exits $[E_{\min}, E_{\max}]$:
    \[
      \mathrm{TSVR} = \frac{1}{T}\sum_{t=1}^{T}
        \mathbf{1}\!\bigl[E_t^{\mathrm{true}} < E_{\min}
          \;\lor\; E_t^{\mathrm{true}} > E_{\max}\bigr].
    \]
    The primary safety metric.  Any controller with $\mathrm{TSVR} > 0$
    has experienced at least one real SOC excursion.

  \item[Severity (P95).]
    The 95th percentile of violation magnitude across violating steps:
    \[
      \mathrm{Sev}_{95} = Q_{0.95}\!\Bigl(\bigl\{
        |E_t^{\mathrm{true}} - E_{\mathrm{nearest\_bound}}|
        : E_t^{\mathrm{true}} \notin [E_{\min}, E_{\max}]\bigr\}\Bigr).
    \]
    Measured in MWh, severity captures how far outside the safe envelope
    the battery ventured.

  \item[IR (Intervention Rate).]
    Fraction of steps where the shield modified the candidate action:
    \[
      \mathrm{IR} = \frac{1}{T}\sum_{t=1}^{T}
        \mathbf{1}\!\bigl[a_t^{\mathrm{safe}} \neq a_t^\star\bigr].
    \]
\end{description}

\subsection{Economic Metrics}

\begin{description}
  \item[Useful work (MWh).]
    Total energy throughput that generated non-negative revenue.
  \item[Expected cost (USD).]
    Dispatch cost under the tariff schedule, including degradation penalty.
  \item[Cost delta (\%).]
    Percentage change in cost relative to the deterministic-LP baseline.
\end{description}

\subsection{Environmental Metrics}

\begin{description}
  \item[Carbon (kg CO$_2$).]
    Grid carbon displaced by battery discharge, computed using the
    marginal emission factor from the generation-mix model.
\end{description}

% ----------------------------------------------------------------
\section{Benchmark Invariants}

The following invariants are checked programmatically after each episode
and recorded in the certificate stream:

\begin{enumerate}
  \item \textbf{Energy conservation.}  The truth-trajectory SOC at
    step~$t+1$ must equal the SOC at step~$t$ plus the energy balance
    from the executed action, to within floating-point tolerance
    ($\epsilon < 10^{-6}$\,MWh).
  \item \textbf{Power bounds.}  The executed power must satisfy
    $0 \le P_t^{+} \le P_{\max}^{+}$ and
    $0 \le P_t^{-} \le P_{\max}^{-}$ at every step.
  \item \textbf{Mutual exclusion.}  $P_t^{+} \cdot P_t^{-} = 0$ at every
    step (no simultaneous charge and discharge).
  \item \textbf{Certificate completeness.}  Every step must emit a
    certificate with zero missing fields (tracked as
    \texttt{certificate\_missing\_fields} in the results CSV).
  \item \textbf{Determinism.}  Re-running the same (scenario, seed,
    controller) triple must produce bit-identical results.
\end{enumerate}

% ----------------------------------------------------------------
\section{Battery Track as Benchmark Foundation}

The CPSBench Battery Track serves as the empirical foundation for all
subsequent evaluation chapters in this dissertation:

\begin{itemize}
  \item \textbf{Chapter~\ref{ch:cert-half-life-blackout}} uses the
    \texttt{drift\_combo} scenario to measure certificate half-life under
    communication blackout.
  \item \textbf{Chapter~\ref{ch:orius-bench-battery}} extends the battery
    track into a public benchmark release with standardised controller
    interfaces and reproduction instructions.
  \item The locked sweep results in
    \texttt{reports/publication/cpsbench\_merged\_sweep.csv} contain
    the full four-dimensional severity sweep (dropout rate $\times$
    delay $\times$ reorder rate $\times$ spike~$\sigma$) across all
    eight controllers and five seeds.
  \item The fault breakdown in
    \texttt{reports/publication/dc3s\_fault\_breakdown.csv} isolates
    the per-fault-type violation and intervention rates for each
    controller.
\end{itemize}

The battery track is intentionally \emph{single-asset}: it evaluates one
battery on one dataset surface at a time.  Fleet-level and cross-region
extensions are addressed in Chapters~\ref{ch:orius-bench-battery}
and~\ref{ch:cert-half-life-blackout}.

% ----------------------------------------------------------------
\section{Chapter Summary}

This chapter introduced the CPSBench Battery Track, a closed-loop
evaluation harness that couples forecasts to battery dispatch under
replayable telemetry faults.  The dual-trajectory architecture separates
truth state from observed state, enabling safety metrics (TSVR, severity)
to be computed on the actual plant trajectory rather than the controller's
degraded view.  Six fault types---dropout, delay jitter, spike, stale
sensor, drift, and reorder---are swept across configurable severity levels
with deterministic episode seeds.  Three metric families (safety, economic,
environmental) provide a comprehensive evaluation surface.  Benchmark
invariants (energy conservation, power bounds, mutual exclusion, certificate
completeness, determinism) ensure result integrity.  The battery track
forms the empirical foundation for all subsequent evaluation and extension
chapters.
